%!TeX root=../emmatop.tex
\chapter[Chapter \thechapter]{}
\lettrine[lines=4,lraise=0.3]{E}{mma}'s pensive meditations, as she walked home, were not interrupted; but on entering the parlour, she found those who must rouse her. Mr Knightley and Harriet had arrived during her absence, and were sitting with her father.—Mr Knightley immediately got up, and in a manner decidedly graver than usual, said,

»I would not go away without seeing you, but I have no time to spare, and therefore must now be gone directly. I am going to London, to spend a few days with John and Isabella. Have you any thing to send or say, besides the »love,« which nobody carries?«

»Nothing at all. But is not this a sudden scheme?«

»Yes—rather—I have been thinking of it some little time.«

Emma was sure he had not forgiven her; he looked unlike himself. Time, however, she thought, would tell him that they ought to be friends again. While he stood, as if meaning to go, but not going—her father began his inquiries.

»Well, my dear, and did you get there safely?—And how did you find my worthy old friend and her daughter?—I dare say they must have been very much obliged to you for coming. Dear Emma has been to call on Mrs and Miss Bates, Mr Knightley, as I told you before. She is always so attentive to them!«

Emma's colour was heightened by this unjust praise; and with a smile, and shake of the head, which spoke much, she looked at Mr Knightley.—It seemed as if there were an instantaneous impression in her favour, as if his eyes received the truth from hers, and all that had passed of good in her feelings were at once caught and honoured.— He looked at her with a glow of regard. She was warmly gratified—and in another moment still more so, by a little movement of more than common friendliness on his part.—He took her hand;—whether she had not herself made the first motion, she could not say—she might, perhaps, have rather offered it—but he took her hand, pressed it, and certainly was on the point of carrying it to his lips—when, from some fancy or other, he suddenly let it go.—Why he should feel such a scruple, why he should change his mind when it was all but done, she could not perceive.—He would have judged better, she thought, if he had not stopped.—The intention, however, was indubitable; and whether it was that his manners had in general so little gallantry, or however else it happened, but she thought nothing became him more.—It was with him, of so simple, yet so dignified a nature.—She could not but recall the attempt with great satisfaction. It spoke such perfect amity.—He left them immediately afterwards—gone in a moment. He always moved with the alertness of a mind which could neither be undecided nor dilatory, but now he seemed more sudden than usual in his disappearance.

Emma could not regret her having gone to Miss Bates, but she wished she had left her ten minutes earlier;—it would have been a great pleasure to talk over Jane Fairfax's situation with Mr Knightley.—Neither would she regret that he should be going to Brunswick Square, for she knew how much his visit would be enjoyed—but it might have happened at a better time—and to have had longer notice of it, would have been pleasanter.—They parted thorough friends, however; she could not be deceived as to the meaning of his countenance, and his unfinished gallantry;—it was all done to assure her that she had fully recovered his good opinion.—He had been sitting with them half an hour, she found. It was a pity that she had not come back earlier!

In the hope of diverting her father's thoughts from the disagreeableness of Mr Knightley's going to London; and going so suddenly; and going on horseback, which she knew would be all very bad; Emma communicated her news of Jane Fairfax, and her dependence on the effect was justified; it supplied a very useful check,—interested, without disturbing him. He had long made up his mind to Jane Fairfax's going out as governess, and could talk of it cheerfully, but Mr Knightley's going to London had been an unexpected blow.

»I am very glad, indeed, my dear, to hear she is to be so comfortably settled. Mrs Elton is very good-natured and agreeable, and I dare say her acquaintance are just what they ought to be. I hope it is a dry situation, and that her health will be taken good care of. It ought to be a first object, as I am sure poor Miss Taylor's always was with me. You know, my dear, she is going to be to this new lady what Miss Taylor was to us. And I hope she will be better off in one respect, and not be induced to go away after it has been her home so long.«

The following day brought news from Richmond to throw every thing else into the background. An express arrived at Randalls to announce the death of Mrs Churchill! Though her nephew had had no particular reason to hasten back on her account, she had not lived above six-and-thirty hours after his return. A sudden seizure of a different nature from any thing foreboded by her general state, had carried her off after a short struggle. The great Mrs Churchill was no more.

It was felt as such things must be felt. Every body had a degree of gravity and sorrow; tenderness towards the departed, solicitude for the surviving friends; and, in a reasonable time, curiosity to know where she would be buried. Goldsmith tells us, that when lovely woman stoops to folly, she has nothing to do but to die; and when she stoops to be disagreeable, it is equally to be recommended as a clearer of ill-fame. Mrs Churchill, after being disliked at least twenty-five years, was now spoken of with compassionate allowances. In one point she was fully justified. She had never been admitted before to be seriously ill. The event acquitted her of all the fancifulness, and all the selfishness of imaginary complaints.

»Poor Mrs Churchill! no doubt she had been suffering a great deal: more than any body had ever supposed—and continual pain would try the temper. It was a sad event—a great shock—with all her faults, what would Mr Churchill do without her? Mr Churchill's loss would be dreadful indeed. Mr Churchill would never get over it.«—Even Mr Weston shook his head, and looked solemn, and said, »Ah! poor woman, who would have thought it!« and resolved, that his mourning should be as handsome as possible; and his wife sat sighing and moralising over her broad hems with a commiseration and good sense, true and steady. How it would affect Frank was among the earliest thoughts of both. It was also a very early speculation with Emma. The character of Mrs Churchill, the grief of her husband—her mind glanced over them both with awe and compassion—and then rested with lightened feelings on how Frank might be affected by the event, how benefited, how freed. She saw in a moment all the possible good. Now, an attachment to Harriet Smith would have nothing to encounter. Mr Churchill, independent of his wife, was feared by nobody; an easy, guidable man, to be persuaded into any thing by his nephew. All that remained to be wished was, that the nephew should form the attachment, as, with all her goodwill in the cause, Emma could feel no certainty of its being already formed.

Harriet behaved extremely well on the occasion, with great self-command. What ever she might feel of brighter hope, she betrayed nothing. Emma was gratified, to observe such a proof in her of strengthened character, and refrained from any allusion that might endanger its maintenance. They spoke, therefore, of Mrs Churchill's death with mutual forbearance.

Short letters from Frank were received at Randalls, communicating all that was immediately important of their state and plans. Mr Churchill was better than could be expected; and their first removal, on the departure of the funeral for Yorkshire, was to be to the house of a very old friend in Windsor, to whom Mr Churchill had been promising a visit the last ten years. At present, there was nothing to be done for Harriet; good wishes for the future were all that could yet be possible on Emma's side.

It was a more pressing concern to shew attention to Jane Fairfax, whose prospects were closing, while Harriet's opened, and whose engagements now allowed of no delay in any one at Highbury, who wished to shew her kindness—and with Emma it was grown into a first wish. She had scarcely a stronger regret than for her past coldness; and the person, whom she had been so many months neglecting, was now the very one on whom she would have lavished every distinction of regard or sympathy. She wanted to be of use to her; wanted to shew a value for her society, and testify respect and consideration. She resolved to prevail on her to spend a day at Hartfield. A note was written to urge it. The invitation was refused, and by a verbal message. »Miss Fairfax was not well enough to write;« and when Mr Perry called at Hartfield, the same morning, it appeared that she was so much indisposed as to have been visited, though against her own consent, by himself, and that she was suffering under severe headaches, and a nervous fever to a degree, which made him doubt the possibility of her going to Mrs Smallridge's at the time proposed. Her health seemed for the moment completely deranged—appetite quite gone—and though there were no absolutely alarming symptoms, nothing touching the pulmonary complaint, which was the standing apprehension of the family, Mr Perry was uneasy about her. He thought she had undertaken more than she was equal to, and that she felt it so herself, though she would not own it. Her spirits seemed overcome. Her present home, he could not but observe, was unfavourable to a nervous disorder:—confined always to one room;—he could have wished it otherwise—and her good aunt, though his very old friend, he must acknowledge to be not the best companion for an invalid of that description. Her care and attention could not be questioned; they were, in fact, only too great. He very much feared that Miss Fairfax derived more evil than good from them. Emma listened with the warmest concern; grieved for her more and more, and looked around eager to discover some way of being useful. To take her—be it only an hour or two—from her aunt, to give her change of air and scene, and quiet rational conversation, even for an hour or two, might do her good; and the following morning she wrote again to say, in the most feeling language she could command, that she would call for her in the carriage at any hour that Jane would name—mentioning that she had Mr Perry's decided opinion, in favour of such exercise for his patient. The answer was only in this short note:

»Miss Fairfax's compliments and thanks, but is quite unequal to any exercise.«

\begin{figure}[tbph]
\centering
\includegraphics[width=.9\linewidth]{45carriage}
\caption{Miss Bates came to the carriage door}
\end{figure}

Emma felt that her own note had deserved something better; but it was impossible to quarrel with words, whose tremulous inequality shewed indisposition so plainly, and she thought only of how she might best counteract this unwillingness to be seen or assisted. In spite of the answer, therefore, she ordered the carriage, and drove to Mrs Bates's, in the hope that Jane would be induced to join her—but it would not do;—Miss Bates came to the carriage door, all gratitude, and agreeing with her most earnestly in thinking an airing might be of the greatest service—and every thing that message could do was tried—but all in vain. Miss Bates was obliged to return without success; Jane was quite unpersuadable; the mere proposal of going out seemed to make her worse.—Emma wished she could have seen her, and tried her own powers; but, almost before she could hint the wish, Miss Bates made it appear that she had promised her niece on no account to let Miss Woodhouse in. »Indeed, the truth was, that poor dear Jane could not bear to see any body—any body at all—Mrs Elton, indeed, could not be denied—and Mrs Cole had made such a point—and Mrs Perry had said so much—but, except them, Jane would really see nobody.«

Emma did not want to be classed with the Mrs Eltons, the Mrs Perrys, and the Mrs Coles, who would force themselves anywhere; neither could she feel any right of preference herself—she submitted, therefore, and only questioned Miss Bates farther as to her niece's appetite and diet, which she longed to be able to assist. On that subject poor Miss Bates was very unhappy, and very communicative; Jane would hardly eat any thing:—Mr Perry recommended nourishing food; but every thing they could command (and never had any body such good neighbours) was distasteful.

Emma, on reaching home, called the housekeeper directly, to an examination of her stores; and some arrowroot of very superior quality was speedily despatched to Miss Bates with a most friendly note. In half an hour the arrowroot was returned, with a thousand thanks from Miss Bates, but »dear Jane would not be satisfied without its being sent back; it was a thing she could not take—and, moreover, she insisted on her saying, that she was not at all in want of any thing.«

When Emma afterwards heard that Jane Fairfax had been seen wandering about the meadows, at some distance from Highbury, on the afternoon of the very day on which she had, under the plea of being unequal to any exercise, so peremptorily refused to go out with her in the carriage, she could have no doubt—putting every thing together—that Jane was resolved to receive no kindness from her. She was sorry, very sorry. Her heart was grieved for a state which seemed but the more pitiable from this sort of irritation of spirits, inconsistency of action, and inequality of powers; and it mortified her that she was given so little credit for proper feeling, or esteemed so little worthy as a friend: but she had the consolation of knowing that her intentions were good, and of being able to say to herself, that could Mr Knightley have been privy to all her attempts of assisting Jane Fairfax, could he even have seen into her heart, he would not, on this occasion, have found any thing to reprove.